\documentclass[12pt]{article}
\usepackage{fontspec}
\usepackage{polyglossia}
\setmainlanguage{french}
\usepackage{csquotes}

\usepackage{tikz}% faire des schémas
\usepackage{pgfplots}% graphiques à partir de données (courbes, histogrammes)
\usepackage{pgfplotstable}% importer les données depuis un fichier csv
\pgfplotsset{compat=1.16}
\pgfplotsset{/pgf/number format/1000 sep=}% pas de séparateur pour les nombres au-dessus de mille
%\pgfplotsset{/pgf/number format/read comma as period}% la virgule est lue comme séparateur décimal (notation française)
%\pgfplotstableset{%
%	assign column name/.style={/pgfplots/table/column name={\textbf{#1}}},% le titre des colonnes apparaît en gras
%	column type=l}% et les colonnes sont alignées à gauche

% Fichiers nécessaires (dans le même dossier que ce fichier) :
% travail_domestique_insee.csv, classement_prenoms_insee.csv, insee_categories_socio_professionnelles.ods

% Dans chaque partie, décommentez le code proposé et complétez les accolades vides.

\begin{document}

\section{Du fichier CSV au tableau}

% 1. Stockez les données de travail_domestique_insee.csv dans une variable, appelée par exemple \travail,
%    puis affichez-les sous forme de tableau.

%\pgfplotstableread[col sep=comma]{}{}

%\pgfplotstabletypeset{}


\section{Du fichier CSV au graphique}

% 2. Affichez les mêmes données sous forme de diagramme : une courbe pour les hommes, une pour les femmes.
%    Dans les crochets de \addplot table, indiquez les colonnes à utiliser : x=..., y=...
%    Dans les accolades, indiquez la variable qui contient les données.

%\begin{tikzpicture}
%	\begin{axis}[xtick=data]
%		\addplot table[]{};
%		\addplot table[]{};
%		\legend{}
%	\end{axis}
%\end{tikzpicture}


\section{Des données textuelles}

% 3. Importez classement_prenoms_insee.csv sous forme d'histogramme en barres.
%    Les prénoms ne sont pas des nombres : ils servent d'étiquettes sur l'axe horizontal.
%    - xticklabels from table={variable}{nom de la colonne des étiquettes}
%    - dans \addplot table : x expr=\coordindex, y=nom de la colonne des valeurs

%\pgfplotstableread[col sep=comma]{}{}

%\begin{tikzpicture}
%	\begin{axis}[ybar, xtick=data, xticklabels from table={}{}, x tick label style={rotate=90}]
%		\addplot table[x expr=\coordindex, y=]{};
%	\end{axis}
%\end{tikzpicture}


\section{Nettoyer ses données}

% 4. Ouvrez insee_categories_socio_professionnelles.ods dans un tableur :
%    - supprimez tout ce qui n'est pas le tableau lui-même (titre, unité, notes, source) ;
%    - enregistrez au format CSV, avec la tabulation comme séparateur (les données contiennent des virgules).
%    Importez ensuite le fichier sous forme de tableau.
%    Décommentez dans le préambule la ligne « read comma as period » ;
%    la première colonne contient du texte : il faut le préciser avec l'option string type.

%\pgfplotstableread[col sep=tab]{}{}

%\pgfplotstabletypeset[columns/Catégorie socioprofessionnelle/.style={string type}]{}

% 5. Décommentez dans le préambule les lignes \pgfplotstableset et observez vos tableaux.

\end{document}
